\documentclass{situationist-preprint}

\title{The Exclusion Argument Excludes Too Much}
\msid{MND-11-b}
\msfamily{MND-11}
\msdiscipline{Philosophy of Mind}
\mstype{Companion manuscript}
\mskeywords{mental causation; exclusion argument; generalization argument; levels; proportionality}

\begin{document}
\maketitle

\begin{abstract}
Kim's exclusion argument is meant to show that irreducible mental properties are causally idle. We argue that its exclusion principle, in the form the argument needs, excludes causes at every level above the fundamental: the causal powers of molecular, cellular and neural properties drain away exactly as those of mental properties do, and if there is no fundamental level, no property causes anything. Kim's attempt to block the generalization by distinguishing levels from orders fails, because the premise that does the work concerns realization, not mereology. The proportionality account of MND-11-a replaces the principle with one that excludes only overly specific causes, and so leaves causes in place at every level at which an effect is proportional to them.
\end{abstract}

\section{The generalization argument}

The exclusion argument has three premises \citep{Kim1998,Kim2005}. \emph{Closure}: every physical effect that has a cause at time $t$ has a sufficient physical cause at $t$. \emph{Exclusion}: no effect has more than one sufficient cause at a time, unless it is a case of genuine overdetermination. \emph{Non-overdetermination}: the effects of mental properties are not systematically overdetermined. If a mental property $M$ is distinct from its physical realizer $P$, and $P$ is sufficient for the effect, Exclusion awards the effect to $P$.

Nothing in this reasoning depends on $M$'s being mental. What it uses is that $M$ is realized by a distinct property $P$ that is sufficient for the effect. Let a \emph{realization chain} for an effect be a sequence of properties $Q_1, Q_2, \dots$ such that each $Q_{n+1}$ realizes $Q_n$ and each is sufficient for the effect.

\begin{proposition}\label{p:gen}
If Exclusion and Non-overdetermination hold for realized and realizing properties, then in any realization chain for an effect, no property that has a realizer in the chain is a cause of the effect.
\end{proposition}

\begin{proof}
Let $Q_n$ have the realizer $Q_{n+1}$ in the chain. Both are sufficient for the effect, and they are distinct. Since the effect is not overdetermined by a property and its own realizer, Exclusion awards it to at most one. The argument that awards it to $Q_{n+1}$ in the mental case is that the realizer is the more basic property and its causal sufficiency is guaranteed by Closure at its level; the same reason applies to any $Q_{n+1}$ and $Q_n$. So $Q_n$ is excluded.
\end{proof}

Take Jones's reaching, caused, we ordinarily say, by her decision to drink. The decision is realized by a neural state; the neural state by a pattern of membrane potentials; that pattern by ion concentrations; those by molecular configurations. Each is sufficient for the reaching, and by Proposition~\ref{p:gen} each is excluded by the next. The neural state is as idle as the decision. This is the drainage worry of \citet{Block2003}: causal power drains down the realization chain to whatever lies at its end.

\begin{corollary}\label{c:drain}
If every realization chain ends in a fundamental property, only fundamental properties are causes. If some chain does not end, then for the effects that chain is sufficient for, no property in it is a cause.
\end{corollary}

The first disjunct makes biology, chemistry and neuroscience causally idle, which no one who accepts Closure in the first place believes. The second is worse, since there is no empirical assurance that physics has a bottom level \citep{Schaffer2003}; if it does not, the argument concludes that nothing causes anything. A premise with this consequence is not a premise to accept on the strength of the mental case.

\section{Kim's reply}

\begin{objection}[Levels and orders]
\citet{Kim1998} distinguishes \emph{levels}, defined by mereology, from \emph{orders}, defined by realization. A second-order property, such as having some property or other that plays a role, is realized by a first-order property at the same level. Exclusion operates only between orders. A micro-based property of a whole, being made of parts with certain properties so related, belongs to the whole and is a first-order property at the higher level. So neural and molecular properties are not excluded by their parts' properties, and the generalization fails.
\end{objection}

\begin{reply}
The reply restores neural causation only by making the relevant neural properties micro-based, that is, by identifying them with configurations of their parts' properties. But the neural properties that figure in the explanations of neuroscience are themselves functionally individuated and multiply realized: a firing pattern can be implemented by many configurations of channels and molecules \citep{Block2003}. These properties stand to their realizers exactly as the mental stands to the neural, and the levels-and-orders distinction puts them on the side of the excluded. Kim's reply therefore saves only properties that are identical to particular micro-configurations, which are just the overly specific properties that do not figure in the generalizations of the special sciences. The generalization argument concerns realization, and Kim's restriction does not touch realization.
\end{reply}

\begin{objection}[Bite the bullet]
One might accept that only fundamental properties are causes, and treat higher-level causal talk as a useful approximation.
\end{objection}

\begin{reply}
The concession costs the exclusionist the reason for the argument. Its point was that the mental is in worse standing than the neural; if the neural is idle too, mental properties are no worse off than any property that science investigates, and the argument yields no special conclusion about the mind \citep{Bennett2003}. It also leaves the conclusion hostage to the existence of a fundamental level, by Corollary~\ref{c:drain}.
\end{reply}

\section{What replaces Exclusion}

The proportionality account of MND-11-a shows why the generalization arises and how to avoid it. Exclusion treats any two sufficient causes as competitors. Principle P1 of MND-11-a awards the effect to the cause that is proportional to it (Definition~2.2), and Proposition~3.3 shows that the proportional cause is unique up to coextension. The correct exclusion principle excludes the overly specific property (Definition~2.3), not the realized one, so that Theorem~3.4 of MND-11-a keeps the mental cause and excludes its maximally specific realizer \citep{Yablo1992,ListMenzies2009}. Applied along a realization chain, the account does not drain: at each level, what is excluded is the property that is too specific for the effect, and what remains is the property whose presence and absence the effect tracks. Which property that is depends on the effect, not on the level. A molecular property can be the proportional cause of a molecular effect, a neural property of a neural one, and a mental property of an action.

\section{Conclusion}

The exclusion principle, as the exclusion argument needs it, makes every realized property idle, and so either confines causation to the fundamental level or abolishes it. Kim's restriction to orders does not help, since the properties of the special sciences are realized. The fault lies with the principle, and the proportionality account of MND-11-a corrects it: mental causes are proportional causes, and physical realizers are overly specific.

\begin{thebibliography}{9}\small
\bibitem[Bennett(2003)]{Bennett2003} Bennett, K. (2003). Why the exclusion problem seems intractable, and how, just maybe, to tract it. \emph{No\^us}, 37(3), 471--497.
\bibitem[Block(2003)]{Block2003} Block, N. (2003). Do causal powers drain away? \emph{Philosophy and Phenomenological Research}, 67(1), 133--150.
\bibitem[Kim(1998)]{Kim1998} Kim, J. (1998). \emph{Mind in a Physical World}. MIT Press.
\bibitem[Kim(2005)]{Kim2005} Kim, J. (2005). \emph{Physicalism, or Something Near Enough}. Princeton University Press.
\bibitem[List and Menzies(2009)]{ListMenzies2009} List, C. and Menzies, P. (2009). Nonreductive physicalism and the limits of the exclusion principle. \emph{The Journal of Philosophy}, 106(9), 475--502.
\bibitem[Schaffer(2003)]{Schaffer2003} Schaffer, J. (2003). Is there a fundamental level? \emph{No\^us}, 37(3), 498--517.
\bibitem[Yablo(1992)]{Yablo1992} Yablo, S. (1992). Mental causation. \emph{The Philosophical Review}, 101(2), 245--280.
\end{thebibliography}

\end{document}
